% Recovered verbatim from chapters_21_25_with_part/chapters_21_25_with_part.tex
% Canonical chapter 21
\chapter{Sub-Planckian Phase Order and Gradient Dynamics}

\section{Phase Order Below the Planck Scale}

In MK--43 Ultra, the fundamental phase field $\tau(x)$ admits sub-Planckian structure described by gradients
\[
\nabla_\mu \tau,\qquad \nabla_\mu \nabla_\nu \tau,
\]
whose fluctuations encode local time-shell deformation.
These gradients are well-defined even in regimes where classical spacetime is not, enabling the model to describe ordering far beneath $\ell_P$.

\section{Gradient Dynamics and Stability}

Define the phase-gradient magnitude
\[
\Phi^2 = g^{\mu\nu} \nabla_\mu \tau\, \nabla_\nu \tau.
\]
Sub-Planckian ordering stabilises when the local cost functional
\[
\mathcal{L}_\tau = \alpha\, \Phi^2 + \beta\, (\Box \tau)^2 + \gamma\, (\nabla \tau)^4
\]
is minimised.
The coefficients $\alpha,\beta,\gamma$ are emergent from ER-foam couplings, ensuring that phase order inherits microstructure from entanglement patterns.

